\documentclass[11pt]{article}

\usepackage{amsmath,amssymb}
\usepackage{xurl}
\usepackage[hidelinks]{hyperref}
\setlength{\emergencystretch}{2em}

\input{command}

\title{Naimark dilation: projective-submeasurement form and the full bipartite assembly}
\author{}
\date{2026-05-08}

\begin{document}
\maketitle

\section{At a glance}

\begin{itemize}
  \item \textbf{Difficulty: low.}  The mathematical issue is now resolved in
  Lean; the note remains to record the faithful projective-submeasurement
  reading of the paper statement.
  \item \textbf{Estimated weight:} no remaining proof construction is expected
  for Naimark itself.
  \item \textbf{Mathlib/project split:} approximately \textbf{60\% Mathlib} and
  \textbf{40\% project-local}.  The proof uses finite-dimensional matrix and
  trace algebra from Mathlib, with project-local packaging for measurements and
  bipartite tensor placement.
  \item \textbf{Status: formalized in the projective-submeasurement form.}  The
  one-measurement Naimark step, the restriction to the original outcomes, and
  the two-sided tensor-product trace identity are now proved in Lean.
  \item \textbf{Remaining note:} the paper's word ``measurement'' must be read
  in the projective-submeasurement sense supplied by the helper lemma.  The
  complete-measurement form on the original outcome type is false for arbitrary
  submeasurements because the residual mass is carried by the additional
  \(\bot\) outcome.
  \item \textbf{Key Mathlib inputs:} finite-dimensional matrix trace
  identities, Naimark's one-measurement compression identity, tensor-product
  states, and the register permutation used by the product-extension state.
\end{itemize}

\section{Key theorem forms}

The note distinguishes the following theorem forms.
\begin{enumerate}
  \item \textbf{Source theorem form: Naimark dilation
  (\texttt{thm:naimark}, displayed at
  \path{references/ldt-paper/orthonormalization.tex}, lines~36--115).}
  For every collection of bipartite
  submeasurements $\{A^x\}_{x\in Q_A}$, $\{B^y\}_{y\in Q_B}$ on a finite
  Hilbert space~$\mathcal H$, there is an enlarged Hilbert space $\widetilde
  {\mathcal H}\supseteq\mathcal H$, an isometry $V:\mathcal H\to\widetilde
  {\mathcal H}$, and projective submeasurements $\{\widetilde A^x\}$,
  $\{\widetilde B^y\}$ on $\widetilde{\mathcal H}\otimes\widetilde{\mathcal H}$
  with
  \[
    \langle\psi|A^x_a\otimes B^y_b|\psi\rangle
    =\langle V\psi|\widetilde A^x_a\otimes \widetilde B^y_b|V\psi\rangle.
    \tag{\texttt{thm:naimark}}
  \]
  This is the displayed paper statement at
  \path{references/ldt-paper/orthonormalization.tex}, lines~36--115, read
  through the helper lemma at lines~121--159.  The helper produces a
  projective submeasurement on the original outcomes; the residual mass is the
  additional $\bot$ outcome in the auxiliary register.
  \item \textbf{Formalized form: one-measurement Naimark.}  For a single
  submeasurement $A$ with outcomes $\mathcal A$ on $\mathcal H$, there is an
  isometric embedding $\mathcal H\to\mathcal H\otimes\C^{\mathrm{Option}(\mathcal A)}$
  and a projective measurement $\widetilde A$ indexed by
  $\mathrm{Option}(\mathcal A)$ whose compression of the outcome
  $(a,\mathrm{some})$ is $A_a$.  The declaration
  \leanid{OneMeasNaimarkData.toProjSubMeas} restricts the completed
  $\mathrm{Option}$-outcome measurement to a projective submeasurement on the
  original outcomes, and
  \leanid{OneMeasNaimarkData.twoSidedCorrelationPreservation} records the
  four-register trace identity for a pair of such dilations.%
  \footnote{See
  \doclink{MIPStarRE/LDT/MakingMeasurementsProjective/NaimarkOneMeas.lean}.}
  \item \textbf{Downstream use.}  The simplified proof of the main theorem
  dilates the slice measurements produced by self-improvement simultaneously,
  with one common ancilla and embedding for all slices, and compresses the
  pasted measurement back to the original space.  Only compression of
  individual effects is used, so neither the bipartite theorem form nor
  preservation of the state-dependent distance is needed.
\end{enumerate}

\section{The source assertion}

We examine \cite{Ji2020LowIndividualDegree}, the Naimark dilation theorem
labelled \texttt{thm:naimark}.%
\footnote{Tracked alongside ledger~\ghissue{1379}.  The paper passages are
\path{references/ldt-paper/orthonormalization.tex}, lines~36--115 (the displayed
theorem statement), 121--159 (the helper lemma \texttt{lem:naimark-helper} and
its proof), 161--187 (the proof of \texttt{thm:naimark} that tensors the
per-question helpers), and 189--272 (the genuine non-preservation example
\texttt{ex:easy-but-long}).}

\medskip

\noindent\emph{Notation.}  We write $X\approx_\delta Y$ on the state $\ket\psi$
to mean the state-dependent operator distance
\(
  \langle\psi|(X-Y)^{\dagger}(X-Y)|\psi\rangle\le\delta,
\)
defined as in \cite{Ji2020LowIndividualDegree} at
\texttt{def:approx\_delta} (\path{references/ldt-paper/preliminaries.tex}, line~348).

\medskip

In the notation of the paper, the source assumes
\[
  \mathcal H=\text{a finite-dimensional Hilbert space},
  \quad \{A^x\}_{x\in Q_A},\,\{B^y\}_{y\in Q_B}\text{ submeasurements on }\mathcal H,
\]
and concludes the bipartite preservation identity above.  The proof proceeds
question by question via \texttt{lem:naimark-helper}.  That helper explicitly
produces a projective submeasurement indexed by the original outcomes.  This is
the mathematically necessary form for arbitrary submeasurements.  If one asks
instead for a complete projective measurement indexed by the original outcomes,
the assertion is false: take a one-outcome submeasurement $A_*=0$ on Alice's
space and a one-outcome measurement $B_*=I$ on Bob's space.  The original
correlation is zero.  But a one-outcome projective measurement on the enlarged
Alice and Bob spaces must have its unique effect equal to the identity, so any
normalized dilated state has correlation one.

The example
\texttt{ex:easy-but-long} demonstrates that the dilation does \emph{not}
preserve the state-dependent distance $\approx_\delta$ defined just above,
which is why downstream arguments in the LDT paper carry the dilation only as
far as the single-outcome marginals.

\section{Comparison with the formalization}

The formal development proves the one-measurement Naimark theorem in the
projective-submeasurement form on the original outcome types, together with
its completed version indexed by \(\mathrm{Option}(\mathsf{Outcome})\).  The
bipartite identity of the source theorem form follows from the four-register
trace identity \leanid{OneMeasNaimarkData.twoSidedCorrelationPreservation}.
The main theorem uses these dilations only through the compression of
individual effects, as described in the simplified proof.

\section{What remains documented}

The former bipartite-assembly gap is closed.  The note remains because it
records a genuine correction to the source wording:
\begin{itemize}
  \item \textbf{Distance non-preservation.}  Example~\texttt{ex:easy-but-long}
  shows that the dilation does not preserve the state-dependent distance
  $\approx_\delta$ recalled in Section~3, so any
  bipartite statement formulated in terms of state-dependent distances would
  need to track an explicit error inflation that is not used downstream.
  \item \textbf{Projective-submeasurement conclusion.}  The helper lemma
  produces projectors on \(\mathrm{Option}(\mathcal A)\).  Restricting to the
  original outcomes gives a projective submeasurement.  It does not give a
  complete projective measurement on the original outcome type unless the
  original submeasurement was already complete.
  \item \textbf{Blueprint interpretation.}  The blueprint therefore states and
  links the theorem in the projective-submeasurement form, with this note
  recording why that is the faithful formal reading of the cited argument.
\end{itemize}

\section{Conclusion}

The Lean declaration \leanid{OneMeasNaimarkData.toProjSubMeas} proves the local
passage from the completed \(\mathrm{Option}\)-outcome projective measurement
to the projective submeasurement on the original outcomes.  The source-facing
statement should be read in the projective-submeasurement form supplied by the
paper's helper lemma; the complete-measurement form on the original outcome
type would be false for submeasurements.

\bibliographystyle{alpha}
\bibliography{references}

\end{document}
